\clearpage
\section*{PHỤ LỤC A. BẢNG ĐỐI CHIẾU KẾT QUẢ CHI TIẾT}
\addcontentsline{toc}{section}{PHỤ LỤC A. BẢNG ĐỐI CHIẾU KẾT QUẢ CHI TIẾT}

Bảng~\ref{tab:detail} dựng lại kết quả chi tiết của từng biến thể từ Bảng I và
Bảng II trong bài báo. Ký hiệu \checkmark{} là không có lỗi, $\times$ là có lỗi.

\begin{table}[H]
\centering
\small
\caption{Kết quả chi tiết theo biến thể (dựng lại từ Bảng I và II của bài báo).}
\label{tab:detail}
\begin{tabular}{@{}llcccc@{}}
\toprule
\textbf{Biến thể} & \textbf{Tiêu chí} & \textbf{ChatGPT-5} & \textbf{ChatGPT-4} & \textbf{Claude} & \textbf{LLaMA} \\
\midrule
\multirow{3}{*}{Coroutine}
 & Build       & \checkmark & \checkmark & \checkmark & $\times$ \\
 & Runtime     & \checkmark & $\times$   & $\times$   & $\times$ \\
 & Correctness & \checkmark & \checkmark & \checkmark & $\times$ \\
\midrule
\multirow{3}{*}{Bất đồng bộ}
 & Build       & \checkmark & \checkmark & \checkmark & \checkmark \\
 & Runtime     & \checkmark & $\times$   & \checkmark & $\times$ \\
 & Correctness & \checkmark & \checkmark & \checkmark & $\times$ \\
\midrule
\multirow{3}{*}{Thuật toán song song}
 & Build       & \checkmark & \checkmark & \checkmark & $\times$ \\
 & Runtime     & \checkmark & $\times$   & $\times$   & $\times$ \\
 & Correctness & \checkmark & \checkmark & \checkmark & \checkmark \\
\midrule
\multirow{3}{*}{OpenMP}
 & Build       & \checkmark & \checkmark & \checkmark & \checkmark \\
 & Runtime     & \checkmark & $\times$   & \checkmark & $\times$ \\
 & Correctness & \checkmark & \checkmark & \checkmark & \checkmark \\
\midrule
\multirow{3}{*}{MPI+OpenMP}
 & Build       & \checkmark & \checkmark & \checkmark & \checkmark \\
 & Runtime     & \checkmark & \checkmark & \checkmark & $\times$ \\
 & Correctness & \checkmark & \checkmark & \checkmark & \checkmark \\
\bottomrule
\end{tabular}
\end{table}

Bảng~\ref{tab:detail} cho thấy rõ hơn Bảng~\ref{tab:success}: ChatGPT-4 thất bại
ở tầng thực thi cho cả bốn biến thể bộ nhớ chia sẻ, dù không có lỗi biên dịch;
Claude chỉ thất bại ở tầng thực thi của coroutine và thuật toán song song (đạt 3/5
biến thể); LLaMA biên dịch thành công 3/5 biến thể (bất đồng bộ, OpenMP và
MPI+OpenMP), đạt tính đúng sau khi sửa ở 3/5 biến thể (thuật toán song song,
OpenMP và MPI+OpenMP), nhưng thất bại ở tầng thực thi cho cả năm biến thể (xem
thêm ghi chú đối chiếu với Bảng~\ref{tab:success} tại Phần~\ref{sec:ketqua}).

\clearpage
\section*{PHỤ LỤC B. KHUNG MPI+OPENMP THAM CHIẾU}
\addcontentsline{toc}{section}{PHỤ LỤC B. KHUNG MPI+OPENMP THAM CHIẾU}

Đoạn mã dưới đây tóm tắt cấu trúc đúng của một chương trình lai MPI+OpenMP cho
bài toán Mandelbrot. Nó không phải mã do LLM sinh, mà là khung tham chiếu để đối
chiếu khi kiểm duyệt. Hai điểm cần chú ý: dùng \code{MPI\_Gatherv} khi số hàng
không chia hết cho số tiến trình, và đặt mốc kết thúc timer \emph{sau} bước gom
dữ liệu nếu muốn đo toàn bộ chi phí.

\begin{lstlisting}[language=C++,caption={Khung MPI+OpenMP với phân vùng theo hàng.},label={lst:mpi},numbers=left]
int rank, size;
MPI_Init(&argc, &argv);
MPI_Comm_rank(MPI_COMM_WORLD, &rank);
MPI_Comm_size(MPI_COMM_WORLD, &size);
int base = height / size;
int local_h = base + (rank < height % size ? 1 : 0);
int y0 = rank * base + std::min(rank, height % size);
std::vector<unsigned char> local((size_t)local_h * width);

MPI_Barrier(MPI_COMM_WORLD);
double t_start = MPI_Wtime();     // bat dau do sau khi dong bo
#pragma omp parallel for schedule(dynamic)
for (int y = 0; y < local_h; ++y)
    for (int x = 0; x < width; ++x)
        local[(size_t)y * width + x] = escape_pixel(x, y0 + y);

std::vector<int> counts(size), displs(size);
for (int r = 0; r < size; ++r) {
    int h = base + (r < height % size ? 1 : 0);
    counts[r] = h * width;
    displs[r] = (r * base + std::min(r, height % size)) * width;
}
std::vector<unsigned char> image;
if (rank == 0) image.resize((size_t)width * height);
MPI_Gatherv(local.data(), (int)local.size(), MPI_UNSIGNED_CHAR,
            image.data(), counts.data(), displs.data(),
            MPI_UNSIGNED_CHAR, 0, MPI_COMM_WORLD);

double t_end = MPI_Wtime();       // ket thuc do sau khi gather (loai I/O)
if (rank == 0) {
    std::cout << "Compute + gather: " << (t_end - t_start) << " s\n";
    write_pgm("mandel.pgm", image, width, height);
}
MPI_Finalize();
\end{lstlisting}

\clearpage
\section*{PHỤ LỤC C. CÂU HỎI THẢO LUẬN}
\addcontentsline{toc}{section}{PHỤ LỤC C. CÂU HỎI THẢO LUẬN}

Các câu hỏi dưới đây có thể dùng để thảo luận nhóm hoặc kiểm tra mức hiểu bài.

\begin{enumerate}
  \item Vì sao một chương trình có thể biên dịch 100\% nhưng chỉ 20\% chạy không
        lỗi? Cho ví dụ cụ thể từ bài báo.
  \item GCC dùng backend nào cho \code{std::execution::par}? Vì sao
        \code{OMP\_NUM\_THREADS} không điều khiển được số luồng trong trường hợp
        này?
  \item \code{done()} của coroutine trả về \code{true} khi nào? Điều gì xảy ra
        nếu \code{final\_suspend} là \code{suspend\_never}?
  \item Phân biệt P4 và P5. Vì sao lỗi định dạng có thể tạo ảnh đen mà không báo
        lỗi?
  \item Vì sao \code{MPI\_Gather} có thể gây segfault thay vì báo lỗi ngay? Khi
        nào cần dùng \code{MPI\_Gatherv}?
  \item Đọc một đường cong mở rộng: làm sao phân biệt “mở rộng tốt” với “luôn
        chạy trên toàn bộ lõi”?
  \item COCOMO có phù hợp để đánh giá chi phí phát triển mã HPC không? Vì sao?
  \item Nếu phải chọn một mô hình cho công việc hằng ngày, bạn sẽ dựa trên tiêu
        chí nào ngoài tỉ lệ biên dịch?
\end{enumerate}
